% =====================================================================
%  Section 10 — Phase 2 scraped sources S1–S8
% =====================================================================

\section{Phase 2 \textbullet\ The Eight Scraped Sources}
\label{sec:phase2-sources}

Every Phase 2 source was chosen because it fills a specific gap in the
story. None is decorative. We document each as: \emph{What it is},
\emph{Why we used it}, \emph{How we scraped/cleaned it}, and
\emph{What it produced}. The original raw artefacts are cached in
\fn{Phase2/RawData/} so the pipeline is reproducible without hitting
the network again.

\subsection{S1 \textbullet\ Transport-for-Cairo GTFS feed}

\textbf{What it is.} The GTFS (General Transit Feed Specification) feed
published by Transport for Cairo. A standard tabular bundle of stops,
routes, trips, shapes, and fare attributes used by transit agencies
worldwide.

\textbf{Why we used it.} GTFS is the single most defensible
representation of the \emph{officially published} formal network — what
a Google-Maps user would see if they tapped ``transit''. We use it to:
(i) compare the formal published reach against Phase 1's measured
informal reach (Q19); (ii) compute formal fare-per-km for the
Q21 economics comparison; (iii) build the formal-vs-informal feature
in the K-Means segmentation (Q24).

\textbf{How we scraped \& cleaned.}
\begin{itemize}
  \item Raw artefact: \fn{RawData/gtfs\_bus\_metro.zip} (1.3 MB) —
        unzipped into \fn{RawData/gtfs\_bus\_metro/}.
  \item Each \fn{*.txt} file is read with pandas; headers normalised;
        rows filtered to the Cairo bounding box; coordinates reprojected
        to EPSG:32636 for downstream spatial work.
\end{itemize}

\textbf{Outputs.} \fn{gtfs\_stops.csv}, \fn{gtfs\_routes.csv},
\fn{gtfs\_trips.csv}, \fn{gtfs\_shapes.csv} (4.9 MB polylines),
\fn{gtfs\_fare\_attributes.csv}.

\subsection{S3 \textbullet\ Wikipedia Cairo Metro master list}

\textbf{What it is.} The English Wikipedia page
\fn{List\_of\_Cairo\_Metro\_stations}: an HTML table with English
name, Arabic name, line code, opening date, and DMS-formatted
coordinates for every operational station.

\textbf{Why we used it.} It is the only source that simultaneously
provides \emph{all} four fields: name, line, opening year, and
geocoordinates. The opening year is critical — every Phase 2 question
that compares pre/post-2014 (Q14, Q15, H2) depends on it. Wikipedia
is more complete and more current than the official ENR list.

\textbf{How we scraped \& cleaned.}
\begin{itemize}
  \item Raw artefact: \fn{RawData/wiki\_metro\_masterlist.html} (196 KB).
  \item BeautifulSoup parses the station-list table.
  \item Coordinates arrive in DMS format (e.g. \strval{30°02'34"N
        31°15'07"E}). A small parser converts to decimal degrees.
  \item Opening phase derived from year thresholds:
        $\le$1996 = Phase~1, 1996--2012 = Phase~2,
        $>$2012 = Phase~3 (the \kw{post-2014} cohort used by every
        downstream comparison).
\end{itemize}

\textbf{Output.} \fn{metro\_stations.csv} — \stat{89 rows}, 100\%
geocoordinates.

\subsection{S4 \textbullet\ Wikipedia Cairo LRT}

\textbf{What it is.} The Wikipedia page \fn{Cairo\_Light\_Rail\_Transit}
with operational and planned-only LRT stations.

\textbf{Why we used it.} The LRT is the most ridership-controversial
piece of post-2014 infrastructure — it serves the New Administrative
Capital and was built ahead of population. Q14, Q22, and the entire
H2 hypothesis depend on having every operational LRT station's location
to compute its catchment.

\textbf{How we scraped \& cleaned.}
\begin{itemize}
  \item Raw artefact: \fn{RawData/wiki\_lrt.html} (117 KB).
  \item The Wikipedia table provides station names but \stat{0\%
        of coordinates on the page} — they are footnoted to external
        databases.
  \item Coordinate backfill is a two-stage rescue: first a per-station
        \stat{Overpass API query} for OSM nodes named after the station
        (rescued 9 of 20); then a Google-Maps fallback (rescued 7
        more). Net: 12 of 20 stations have valid lat/lon for analysis;
        the remaining 8 are still planned-only.
\end{itemize}

\textbf{Output.} \fn{lrt\_stations.csv} (20 rows), with the backfill
documented in the cleaning notebook's
\fn{lrt\_coordinate\_backfill.html} export.

\figitem{phase2/lrt_coordinate_backfill.png}{1.0}{S4 \textbullet\ LRT coordinate backfill audit. Bars show how many stations were rescued by Overpass vs Google Maps vs remained coordinate-less.}{fig:lrt-backfill}

\subsection{S5 \textbullet\ Google Maps BRT scrape}

\textbf{What it is.} The Cairo Ring-Road BRT has no Wikipedia page and
no Wikidata entry; the only structured representation is on Google Maps.

\textbf{Why we used it.} The \kw{H3 hypothesis} (BRT corridor match) is
the single most narratively important contrast point in Phase 2 —
it is the case where state planning aligns with demand, and we cannot
test it without BRT station coordinates.

\textbf{How we scraped \& cleaned.}
\begin{itemize}
  \item Raw artefact: \fn{RawData/gmaps\_brt\_brt\_en.html} (426 KB)
        + \fn{gmaps\_brt\_trd\_ar.html} (494 KB) — playwright-rendered
        search-result snapshots in English and Arabic.
  \item \stat{10 viewport queries} along the Ring Road.
  \item Regex pulls coordinates out of the JavaScript URL fragments
        (\fn{!3d.../!4d...}).
  \item \kw{Arabic transliteration} via \fn{uroman} (franko) — phonetic,
        not semantic. We tried MarianMT first; it semantically translates
        \ar{المرج} (al-Marg) to ``lawn'' instead of phonetically to
        \strval{almrj}, which broke deduplication.
  \item \kw{Three-tier deduplication}: spatial (within 50\,m), name
        (RapidFuzz token-sort $\geq$ 88), and source-cluster.
\end{itemize}

\textbf{Output.} \fn{brt\_stations.csv} — \stat{21 stations}, all with
coordinates. The BRT scrape grew from an original 12 to the current 21
after S5 was rerun against fresh Google Maps tiles.

\figitem{phase2/brt_scrape_diagnostic.png}{1.0}{S5 \textbullet\ BRT scrape diagnostic — counts before and after the three deduplication tiers.}{fig:brt-scrape}

\subsection{S6 \textbullet\ citypopulation.de — Greater Cairo districts}

\textbf{What it is.} \fn{citypopulation.de} publishes administrative
population data for every country across multiple census years. The
\fn{egypt/admin} page is a single HTML table with \stat{408 admin rows}
covering 27 governorates.

\textbf{Why we used it.} Population is the demand denominator of every
Phase 2 question. Without it we cannot:
\begin{itemize}
  \item compute station-per-100k-residents ratios (the H1 metric),
  \item compute 2-km catchment populations (the H2 metric),
  \item identify the fastest-growing districts (Q16),
  \item segment districts in K-Means (Q24).
\end{itemize}

\textbf{How we scraped \& cleaned.}
\begin{itemize}
  \item Raw artefact: \fn{RawData/citypop\_egypt\_admin.html} (282 KB).
  \item HTML table parsing with BeautifulSoup; year columns reshaped
        from wide to long.
  \item Filtered to Greater Cairo (Cairo, Giza, Qalyubia governorates)
        — 68 districts.
  \item \kw{CAGR} (compound annual growth rate) computed per district
        across 2006$\to$2017 census windows.
  \item District centroids fetched via OSM Nominatim — cached in
        \fn{RawData/nominatim\_districts.json} (5.5 KB). The Nominatim
        centroids replace the earlier draft's three-governorate-centroid
        proxy and are what turned the H1 effect size from medium into
        huge (see~\S\ref{sec:phase2-hypotheses}).
\end{itemize}

\textbf{Outputs.}
\begin{itemize}
  \item \fn{districts.csv} (long: district-year rows, 480+ rows).
  \item \fn{districts\_wide.csv} (wide: one row per district, years
        across columns).
\end{itemize}

\figitem{phase2/districts_cagr_distribution.png}{0.95}{S6 \textbullet\ District-level CAGR distribution (2006$\to$2017). The right tail are the New Administrative Capital and the satellite cities; this is the demand context for Q16 and Q22.}{fig:districts-cagr}

\subsection{S7 \textbullet\ Vehicle ownership · Wikipedia + World Bank}

\textbf{What it is.} National-level motorisation context: a Wikipedia
table of vehicles per capita, plus the World Bank API endpoints
\fn{wb\_motor\_vehicles\_per\_1k} and \fn{wb\_passenger\_cars\_per\_1k}.

\textbf{Why we used it.} The Q18b motorisation question and the
secondary K-Means feature both need a per-governorate vehicle
estimate. Direct governorate-level data is not consistently scrapeable;
this two-source fallback gives us a defensible apportionment by
population share.

\textbf{How we scraped \& cleaned.}
\begin{itemize}
  \item Raw artefacts: \fn{RawData/wiki\_vehicles\_per\_capita.html}
        (353 KB), \fn{wb\_motor\_vehicles\_per\_1k.json},
        \fn{wb\_passenger\_cars\_per\_1k.json}.
  \item Greater Cairo apportionment by 2017 population share.
\end{itemize}

\textbf{Output.} \fn{vehicles\_by\_governorate.csv}.

\subsection{S8 \textbullet\ OpenStreetMap (Overpass API)}

\textbf{What it is.} A spatial query against the OpenStreetMap
Overpass API for every transport-tagged node within the Cairo bounding
box.

\textbf{Why we used it.} OSM is the third independent source family in
the project (after Wikipedia and Google Maps), so it functions as a
\kw{cross-verification layer}. It is also the rescue source for the
LRT coordinate backfill (S4).

\textbf{How we scraped \& cleaned.}
\begin{itemize}
  \item Raw artefact: \fn{RawData/osm\_overpass.json} (114 KB) +
        per-station \fn{osm\_lrt\_per\_station.json} (20 KB) for the
        backfill.
  \item Overpass QL filtered on \fn{amenity}, \fn{public\_transport},
        \fn{highway=bus\_stop}, \fn{railway=station}.
  \item Reprojected to EPSG:32636 and saved as GeoJSON.
\end{itemize}

\textbf{Output.} \fn{osm\_features.geojson} (92 KB, 1{,}000+ points).

\figitem{phase2/osm_cross_verification_map.png}{1.0}{S8 \textbullet\ OSM cross-verification map. Each scraped station's coordinate is paired with the nearest OSM node; the discrepancy distance is visualised.}{fig:osm-cross}

\subsection{Source-by-source summary}

\rowcolors{2}{darkpanel}{darkbg}
\begin{table}[H]
\centering\small
\renewcommand{\arraystretch}{1.25}
\begin{tabular}{@{} >{\color{plotorange}\bfseries}l
                    >{\color{bodytext}}l
                    >{\color{plotyellow}\ttfamily}r
                    >{\color{bodytext}}p{4.5cm} @{}}
\toprule
\color{plotblue}\# & \color{plotblue}Origin & \color{plotblue}Rows & \color{plotblue}Method shorthand \\
\midrule
S1 & Transport-for-Cairo GTFS                     & 50k+ & GTFS spec, bbox filter \\
S3 & Wikipedia metro master list                  & 89   & DMS$\to$decimal, year-threshold phase \\
S4 & Wikipedia LRT                                & 20   & 9 Overpass + 7 GMaps backfill \\
S5 & Google Maps BRT                              & 21   & 10 viewport queries, regex aria-labels, uroman, 3-tier dedup \\
S6 & citypopulation.de                            & 408  & HTML table, Nominatim centroids, CAGR \\
S7 & Wikipedia + World Bank                       &  5   & per-capita apportionment \\
S8 & OpenStreetMap Overpass                       & 1k+  & Overpass QL, GeoJSON \\
\bottomrule
\end{tabular}
\caption{The eight Phase 2 sources at a glance. The two ``missing'' rows (S2, the public-transport route GTFS, was merged into S1; the original draft also reserved an S9 slot) reflect the project's notebook history rather than a data gap.}
\end{table}
\rowcolors{0}{}{}

\subsection{Pre-cleaning null audit (Phase 2)}

\figitem{phase2/null_audit_before_after.png}{1.0}{Phase 2 null audit before vs after cleaning. Each bar shows the percentage of nulls per column per source. The post-cleaning bars all sit at zero or are visually negligible.}{fig:p2-null-audit}
